\subsection{When Did Language Appear?}
\label{sec:timeline}
\slidesrc{Intro: 9--12}

\subsubsection{The timeline of humankind}
\slidesrc{Intro: 9--10}

Table~\ref{tab:harari} places our species on the timeline of human evolution. It is
adapted from Yuval Noah Harari's \emph{Sapiens: A Brief History of Humankind} (2014).
The abbreviations are MYA = million years ago, YA = years ago.

\begin{center}
\small
\begin{tabularx}{\linewidth}{@{}l X@{}}
\toprule
\textbf{When} & \textbf{What happened} \\
\midrule
6 MYA            & Last common ``grandmother'' of humans and chimpanzees \\
2.5 MYA          & Evolution of the genus \emph{Homo} in Africa; first stone tools \\
2 MYA            & \textbf{Humans spread from Africa to Eurasia}; evolution of different human
                   species (\emph{Homo rudolfensis} and \emph{Homo ergaster} in Africa,
                   \emph{Homo erectus} in East Asia) \\
500{,}000 YA     & Neanderthals evolve in Europe and the Middle East \\
200{,}000 YA     & \emph{Homo sapiens} evolves in East Africa \\
70{,}000 YA      & \textbf{Sapiens spreads out of Africa} \\
70{,}000--30{,}000 YA & \textbf{Cognitive revolution}: invention of boats, arrows, needles;
                   beginnings of art, religion, commerce, social stratification,
                   \textbf{and language} \\
45{,}000 YA      & Sapiens settles Australia \\
30{,}000 YA      & Extinction of Neanderthals \\
16{,}000 YA      & Sapiens settles America \\
13{,}000 YA      & Extinction of \emph{Homo floresiensis}: \textbf{Homo sapiens is the only
                   surviving human species} \\
12{,}000 YA      & Agricultural revolution, permanent settlements \\
\bottomrule
\end{tabularx}
\captionof{table}{Milestones of human evolution (adapted from Harari). Note how recent
everything after the cognitive revolution is compared with the 6 million years of the
whole story.}
\label{tab:harari}
\end{center}

The course also shows a map of the movement of \emph{Homo sapiens} out of Africa: from
East Africa into the Middle East, then along the coasts of southern Asia to Australia
(about 50{,}000 years ago), into Europe and northern Asia, and finally across the Bering
region into America (about 15{,}000 years ago), while Neanderthals already lived in
Europe and \emph{Homo erectus} in Asia.

\subsubsection{What recent research says}
\slidesrc{Intro: 11}

Harari places the birth of language in the cognitive revolution, about
70{,}000--30{,}000 years ago. A 2025 study reported by MIT News (``When did human language
emerge?'', 14 March 2025) argues for an earlier date:

\begin{itemize}[nosep]
  \item \textbf{135{,}000 years ago}: genomic and geographic data suggest that the
        biological and cognitive \emph{capacity} for language was already fundamentally
        present. The argument is that the earliest split between human populations
        happened about 135{,}000 years ago, and since \emph{all} human populations have
        language, the capacity must predate the split.
  \item \textbf{100{,}000 years ago}: archaeological records show an explosion of
        symbolic behaviours, like art and complex tools, indicating active social
        communication. Language was being \emph{used} socially.
  \item \textbf{Gradual origins}: scientists widely agree that language developed from a
        complex mix of hand gestures, facial expressions and vocalisations over time,
        not in a single step.
\end{itemize}

\begin{nbwarning}[capacity vs.\ use]
The two dates answer different questions. ``The capacity for language existed
135{,}000 years ago'' is about biology (the brain and vocal tract could support
language). ``Symbolic behaviour exploded 100{,}000 years ago'' is about culture (people
were actually communicating with symbols). Do not mix them up when asked ``when did
language appear?''.
\end{nbwarning}

\subsubsection{Did the Neanderthals speak?}
\slidesrc{Intro: 12}

The course leaves this as a question for discussion. There is no definitive answer,
but these are the main pieces of evidence usually brought up:
\begin{itemize}[nosep]
  \item \textbf{for}: Neanderthals had the same variant as ours of the \emph{FOXP2} gene,
        involved in speech and language; a Neanderthal hyoid bone (the bone that
        supports the tongue) found at Kebara, Israel, is almost identical to ours; they
        made complex tools, used fire, buried their dead and probably made simple
        symbolic objects;
  \item \textbf{against / uncertain}: their symbolic record is much poorer than that of
        \emph{Sapiens} of the same period, and soft tissue (the vocal tract, the brain
        itself) does not fossilise, so the anatomy of speech can only be inferred.
\end{itemize}
Many researchers today think Neanderthals had \emph{some} form of language, though
possibly less complex than ours. Note that if language capacity predates 135{,}000
years ago in \emph{Sapiens}, the question of whether it also predates the split between
\emph{Sapiens} and Neanderthals (several hundred thousand years ago) remains open.
